\documentclass{article}
\usepackage{amsmath}
\usepackage{array}
\usepackage[a4paper, margin=1in]{geometry}

\begin{document}
\title{Discrete Mathematics Notes}
\author{Lochana Ranatunga}
\maketitle

\section{Propositions and Logical Language}
\subsection*{Statements/Propositions}

\begin{itemize}
  \item A proposition is a declarative statement that is either true or false, but not both.
  \item Commands, questions do not have a truth value, therefore not considered propositions in mathematical logic.
  \item Truth value of open statements (ex. $x > 3$ , $x^2 + 1 = 0$) depends on the variable.
\end{itemize}

\subsection*{Logical Connectives}
\begin{enumerate}
  \item AND ($\land$)
  \item OR ($\lor$)
  \item NOT ($\neg$)
  \item IF...THEN ($\rightarrow$)
  \item IF AND ONLY IF ($\leftrightarrow$)
\end{enumerate}

\subsubsection*{Conjunction (AND)}
Symbol: $\land$ \hspace{1.5cm}
Form: $P \land Q$ \hspace{1.5cm}
Meaning: Both P and Q are true \\
\begin{table}[h]
\centering
\renewcommand{\arraystretch}{1.01}
\begin{tabular}{|>{\centering\arraybackslash}p{2cm}|>{\centering\arraybackslash}p{2cm}|>{\centering\arraybackslash}p{2cm}|}
\hline
P & Q &  $P \land Q$\\
\hline
T & T & T \\
\hline
T & F & F \\
\hline
F & T & F \\
\hline
F & F & F \\
\hline
\end{tabular}
\caption{Truth Table for AND}
\end{table}

\subsubsection*{Disjunction (OR)}
Symbol: $\lor$ \hspace{1.5cm}
Form: $P \lor Q$ \hspace{1.5cm}
Meaning: At least one of P or Q are true \\
\begin{table}[h]
\centering
\renewcommand{\arraystretch}{1.01}
\begin{tabular}{|>{\centering\arraybackslash}p{2cm}|>{\centering\arraybackslash}p{2cm}|>{\centering\arraybackslash}p{2cm}|}
\hline
P & Q &  $P \lor Q$\\
\hline
T & T & T \\
\hline
T & F & T \\
\hline
F & T & T \\
\hline
F & F & F \\
\hline
\end{tabular}
\caption{Truth Table for OR}
\end{table}\\
OR is inclusive, unless specifically stated that it is exclusive. (Inclusive OR is true, even when both are true)\\

\subsubsection*{Implication (IF...THEN)}
Symbol: $\rightarrow$ \hspace{1.5cm}
Form: $P \rightarrow Q$ \hspace{1.5cm}
Meaning: If P is true, then Q is true \\
\begin{table}[h]
\centering
\renewcommand{\arraystretch}{1.01}
\begin{tabular}{|>{\centering\arraybackslash}p{2cm}|>{\centering\arraybackslash}p{2cm}|>{\centering\arraybackslash}p{2cm}|}
\hline
P & Q &  $P \rightarrow Q$\\
\hline
T & T & T \\
\hline
T & F & F \\
\hline
F & T & T \\
\hline
F & F & T \\
\hline
\end{tabular}
\caption{Truth Table for IF...THEN}
\end{table}\\
In classical logic, if the antecedent (P) is False, the implication is True irrespective to the truth value of the consequent (Q).
Therefore, this is referred to as 'Material implication'.\\

\subsubsection*{Biconditional (IF AND ONLY IF)}
Symbol: $\leftrightarrow$ \hspace{1.5cm}
Form: $P \leftrightarrow Q$ \hspace{1.5cm}
Meaning: P is true, exactly when Q is true \\
'If and only if' is abbreviated as, 'iff'.\\
\begin{table}[h]
\centering
\renewcommand{\arraystretch}{1.01}
\begin{tabular}{|>{\centering\arraybackslash}p{2cm}|>{\centering\arraybackslash}p{2cm}|>{\centering\arraybackslash}p{2cm}|}
\hline
P & Q &  $P \leftrightarrow Q$\\
\hline
T & T & T \\
\hline
T & F & F \\
\hline
F & T & F \\
\hline
F & F & T \\
\hline
\end{tabular}
\caption{Truth Table for IF AND ONLY IF}
\end{table}\\
\textbf{Biconditional in terms of other logical connectives:}\\
($p \leftrightarrow q) \equiv (p \rightarrow q) \land (q \rightarrow p)$\\
$(p \leftrightarrow q) \equiv (p \land q) \lor (\neg p \land \neg q)$\\

\subsubsection*{Negation (NOT)}
Symbol: $\neg$ \hspace{1.5cm}
Form: $\neg P$ \hspace{1.5cm}
Meaning: not P is True. (P is False) \\
\begin{table}[h]
\centering
\renewcommand{\arraystretch}{1.01}
\begin{tabular}{|>{\centering\arraybackslash}p{2cm}|>{\centering\arraybackslash}p{2cm}|}
\hline
P & $\neg P $\\
\hline
T & F\\
\hline
F & T \\
\hline
\end{tabular}
\caption{Truth Table for NOT}
\end{table}\\

\newpage
\subsubsection*{Exclusive OR (XOR)}
Symbol: $\oplus$ \hspace{1.5cm}
Form: $P \oplus Q$ \hspace{1.5cm}
Meaning:  True exactly when one of P or Q is true\\
\begin{table}[h]
\centering
\renewcommand{\arraystretch}{1.01}
\begin{tabular}{|>{\centering\arraybackslash}p{2cm}|>{\centering\arraybackslash}p{2cm}|>{\centering\arraybackslash}p{2cm}|}
\hline
P & Q &  $P \oplus Q$\\
\hline
T & T & F \\
\hline
T & F & T \\
\hline
F & T & T \\
\hline
F & F & F \\
\hline
\end{tabular}
\caption{Truth Table for XOR}
\end{table}\\
\textbf{Biconditional in terms of other logical connectives:}\\
($p \leftrightarrow q) \equiv (p \rightarrow q) \land (q \rightarrow p)$\\
$(p \leftrightarrow q) \equiv (p \land q) \lor (\neg p \land \neg q)$\\

\end{document}